\documentclass[../main.tex]{subfiles}
\begin{document}

\chapter{同期機構}
\lessoninfo{
  video={P3L4},
  textbook={OSTEP \cite{ostep} ch.~28，31；Silberschatz ら \cite{silberschatz2018} ch.~6--7},
  keywords={スピンロック，ビジーウェイト，セマフォ，リーダ・ライタロック，モニタ，バリア，ウェイトフリー同期，RCU，アトミック命令，テストアンドセット，キャッシュ一貫性，書き込み無効化，テストアンドテストアンドセット，キューイングロック},
}

第3章ではミューテックスと条件変数を導入し，第9章では共有メモリを通じて
通信するプロセスもスレッドと同じように同期を必要とすることを示した．この章
では，ミューテックスと条件変数より単純か，あるいは表現力の高い同期機構
---スピンロック，セマフォ，リーダ・ライタロック，モニタ，およびあまり
一般的でないいくつかの機構---を概観する．続いて，そうした機構がすべて最終的
に依拠するもの，すなわちハードウェアが提供するアトミック命令と，それが共有
メモリ型マルチプロセッサのキャッシュとどう関わるかを扱う．

\section{ミューテックスと条件変数の先へ}
\label{sec:l10-limits}

ミューテックスと条件変数は第3章の同期方針を表現できるが，3つの限界があり，
それが他の機構を必要とする理由となる．\source{P3L4，同期についてさらに；
OSTEP ch.~28．}
\begin{description}
  \item[誤りやすい利用] 正しさはすべてプログラマに委ねられる．共有状態への
    アクセスはすべて正しいミューテックスをロックし，どの経路でもそれを解放
    しなければならず，状態の変化はすべて正しい条件変数にシグナルまたは
    ブロードキャストしなければならない．第3章に挙げた落とし穴の大半，例えば
    誤ったミューテックスのアンロックや誤った条件変数へのシグナルは，これら
    の規則に違反するものであり，機構自体はそのどれも検出しない．
  \item[表現力の限界] ミューテックスが強制する方針はちょうど1つ，一度に1つ
    のスレッドという方針だけである．複数のリーダまたは1つのライタを受け
    入れる，待っているライタを優先するといったより豊かな方針は，第3章の
    リーダ・ライタ問題の解における代理変数のように，補助変数と明示的な検査
    で符号化しなければならない．
  \item[低水準の支援への依存] ミューテックス自体が共有状態である．その
    ロック操作は，ミューテックスが空いていることを見つけ，それを所有済みと
    印を付けることを，1つの不可分な手順で行わなければならない．実用的な実装
    はその手順を，ハードウェアが提供するアトミック命令から得る
    （\cref{sec:l10-hw}）．
\end{description}
セマフォ，リーダ・ライタロック，モニタ，および
\cref{sec:l10-semaphore,sec:l10-rwlock,sec:l10-monitor,sec:l10-other}の機構
は，最初の2つの限界に対処する．\cref{sec:l10-spinlock}のスピンロックは
それらすべての土台となる低水準の構成要素であり，3つ目の限界はこの章の後半の
主題である．

\section{スピンロック}
\label{sec:l10-spinlock}

ミューテックスの次に単純な機構がスピンロックである．\source{P3L4，
スピンロック；OSTEP ch.~28．}

\begin{definition}[スピンロック]
  \term[すぴんろっく]{スピンロック}[spinlock]とは，ミューテックスと同様に
  ロック操作とアンロック操作によって相互排除を提供するロックであり，ロックが
  保持されていることを見つけたスレッドがブロックしない点だけが異なる．
  そのスレッドは実行を続け，ロックが空くか自分がプリエンプトされるまで，
  ロックを繰り返し検査する．
\end{definition}

条件が成立するまで CPU を手放すのではなく，ループの中で条件を何度も検査する
ことを\term[びじーうぇいと]{ビジーウェイト}[busy waiting]と呼び，スピン
ロック上でビジーウェイトするスレッドは\emph{スピンする}という．スピンロック
はミューテックスとまったく同じように使う．
\begin{lstlisting}[language=C, numbers=none]
spinlock_lock(s);
// クリティカルセクション
spinlock_unlock(s);
\end{lstlisting}

2つの機構が異なるのは，ロックが他のスレッドに保持されている間にスレッドが
何をするかだけであり，そのコストも異なる（\cref{fig:l10-spin-block}）．
ミューテックスでブロックしたスレッドは待っている間 CPU 時間を消費しないが，
自分から他のスレッドへのコンテキストスイッチと，ミューテックスが解放された
後の起床および元へ戻るコンテキストスイッチの代価を払う．スピンするスレッド
はどちらのコンテキストスイッチも避けられるが，ロックが保持されている間
ずっと自分の CPU を占有する．\margin{POSIX では \texttt{pthread\_spin\_lock} と
\texttt{pthread\_spin\_unlock}．}

\begin{figure}[htbp]
  \centering
  % Blocking on a mutex versus spinning on a spinlock: timelines of two CPUs.
% Usage: % Blocking on a mutex versus spinning on a spinlock: timelines of two CPUs.
% Usage: % Blocking on a mutex versus spinning on a spinlock: timelines of two CPUs.
% Usage: \input{figures/l10-spin-block}   (width ~116 mm)
\begin{tikzpicture}[x=1mm, y=1mm, font=\footnotesize\sffamily,
  lab/.style={font=\scriptsize\sffamily, align=center},
  row/.style={font=\scriptsize\sffamily, anchor=east},
  ttl/.style={font=\footnotesize\sffamily\bfseries, anchor=west},
  >={Stealth[length=1.6mm]}]
  % \ltenseg{x1}{x2}{y}{fill}{text}
  \def\ltenseg#1#2#3#4#5{%
    \draw[fill=#4] (#1,#3-2.5) rectangle (#2,#3+2.5);
    \node[lab] at ({(#1+#2)/2},#3) {#5};}
  % event lines
  \draw[densely dashed, ln-muted] (50,33) -- (50,-12);
  \draw[densely dashed, ln-muted] (76,33) -- (76,-12);
  \node[lab, anchor=south, fill=white, inner sep=1pt] at (50,33)
    {\iflangja{T2 が lock を呼ぶ}{T2 calls lock}};
  \node[lab, anchor=south, fill=white, inner sep=1pt] at (76,33)
    {\iflangja{T1 が unlock する}{T1 unlocks}};
  % (a) mutex
  \node[ttl] at (0,29) {\iflangja{(a) ミューテックス：T2 はブロック}{(a) mutex: T2 blocks}};
  \node[row] at (22,22) {CPU 1};
  \ltenseg{24}{36}{22}{white}{T1}
  \ltenseg{36}{76}{22}{ln-box}{\iflangja{T1：クリティカルセクション}{T1: critical section}}
  \ltenseg{76}{116}{22}{white}{T1}
  \node[row] at (22,15) {CPU 2};
  \ltenseg{24}{50}{15}{white}{T2}
  \ltenseg{50}{55}{15}{black!35}{}
  \ltenseg{55}{76}{15}{white}{\iflangja{T3（別のスレッド）}{T3 (another thread)}}
  \ltenseg{76}{81}{15}{black!35}{}
  \ltenseg{81}{116}{15}{ln-box}{\iflangja{T2：クリティカルセクション}{T2: critical section}}
  % (b) spinlock
  \node[ttl] at (0,6) {\iflangja{(b) スピンロック：T2 はスピン}{(b) spinlock: T2 spins}};
  \node[row] at (22,-1) {CPU 1};
  \ltenseg{24}{36}{-1}{white}{T1}
  \ltenseg{36}{76}{-1}{ln-box}{\iflangja{T1：クリティカルセクション}{T1: critical section}}
  \ltenseg{76}{116}{-1}{white}{T1}
  \node[row] at (22,-8) {CPU 2};
  \ltenseg{24}{50}{-8}{white}{T2}
  \ltenseg{50}{76}{-8}{ln-caution!20}{\iflangja{ビジーウェイト}{busy waiting}}
  \ltenseg{76}{116}{-8}{ln-box}{\iflangja{T2：クリティカルセクション}{T2: critical section}}
  % legend
  \draw[fill=black!35] (24,-18) rectangle (28,-15);
  \node[lab, anchor=west] at (29,-16.5) {\iflangja{コンテキストスイッチ}{context switch}};
  \draw[fill=ln-box] (64,-18) rectangle (68,-15);
  \node[lab, anchor=west] at (69,-16.5) {\iflangja{ロックを保持}{lock held}};
  \draw[->] (92,-16.5) -- node[lab, above] {\iflangja{時間}{time}} (116,-16.5);
\end{tikzpicture}
   (width ~116 mm)
\begin{tikzpicture}[x=1mm, y=1mm, font=\footnotesize\sffamily,
  lab/.style={font=\scriptsize\sffamily, align=center},
  row/.style={font=\scriptsize\sffamily, anchor=east},
  ttl/.style={font=\footnotesize\sffamily\bfseries, anchor=west},
  >={Stealth[length=1.6mm]}]
  % \ltenseg{x1}{x2}{y}{fill}{text}
  \def\ltenseg#1#2#3#4#5{%
    \draw[fill=#4] (#1,#3-2.5) rectangle (#2,#3+2.5);
    \node[lab] at ({(#1+#2)/2},#3) {#5};}
  % event lines
  \draw[densely dashed, ln-muted] (50,33) -- (50,-12);
  \draw[densely dashed, ln-muted] (76,33) -- (76,-12);
  \node[lab, anchor=south, fill=white, inner sep=1pt] at (50,33)
    {\iflangja{T2 が lock を呼ぶ}{T2 calls lock}};
  \node[lab, anchor=south, fill=white, inner sep=1pt] at (76,33)
    {\iflangja{T1 が unlock する}{T1 unlocks}};
  % (a) mutex
  \node[ttl] at (0,29) {\iflangja{(a) ミューテックス：T2 はブロック}{(a) mutex: T2 blocks}};
  \node[row] at (22,22) {CPU 1};
  \ltenseg{24}{36}{22}{white}{T1}
  \ltenseg{36}{76}{22}{ln-box}{\iflangja{T1：クリティカルセクション}{T1: critical section}}
  \ltenseg{76}{116}{22}{white}{T1}
  \node[row] at (22,15) {CPU 2};
  \ltenseg{24}{50}{15}{white}{T2}
  \ltenseg{50}{55}{15}{black!35}{}
  \ltenseg{55}{76}{15}{white}{\iflangja{T3（別のスレッド）}{T3 (another thread)}}
  \ltenseg{76}{81}{15}{black!35}{}
  \ltenseg{81}{116}{15}{ln-box}{\iflangja{T2：クリティカルセクション}{T2: critical section}}
  % (b) spinlock
  \node[ttl] at (0,6) {\iflangja{(b) スピンロック：T2 はスピン}{(b) spinlock: T2 spins}};
  \node[row] at (22,-1) {CPU 1};
  \ltenseg{24}{36}{-1}{white}{T1}
  \ltenseg{36}{76}{-1}{ln-box}{\iflangja{T1：クリティカルセクション}{T1: critical section}}
  \ltenseg{76}{116}{-1}{white}{T1}
  \node[row] at (22,-8) {CPU 2};
  \ltenseg{24}{50}{-8}{white}{T2}
  \ltenseg{50}{76}{-8}{ln-caution!20}{\iflangja{ビジーウェイト}{busy waiting}}
  \ltenseg{76}{116}{-8}{ln-box}{\iflangja{T2：クリティカルセクション}{T2: critical section}}
  % legend
  \draw[fill=black!35] (24,-18) rectangle (28,-15);
  \node[lab, anchor=west] at (29,-16.5) {\iflangja{コンテキストスイッチ}{context switch}};
  \draw[fill=ln-box] (64,-18) rectangle (68,-15);
  \node[lab, anchor=west] at (69,-16.5) {\iflangja{ロックを保持}{lock held}};
  \draw[->] (92,-16.5) -- node[lab, above] {\iflangja{時間}{time}} (116,-16.5);
\end{tikzpicture}
   (width ~116 mm)
\begin{tikzpicture}[x=1mm, y=1mm, font=\footnotesize\sffamily,
  lab/.style={font=\scriptsize\sffamily, align=center},
  row/.style={font=\scriptsize\sffamily, anchor=east},
  ttl/.style={font=\footnotesize\sffamily\bfseries, anchor=west},
  >={Stealth[length=1.6mm]}]
  % \ltenseg{x1}{x2}{y}{fill}{text}
  \def\ltenseg#1#2#3#4#5{%
    \draw[fill=#4] (#1,#3-2.5) rectangle (#2,#3+2.5);
    \node[lab] at ({(#1+#2)/2},#3) {#5};}
  % event lines
  \draw[densely dashed, ln-muted] (50,33) -- (50,-12);
  \draw[densely dashed, ln-muted] (76,33) -- (76,-12);
  \node[lab, anchor=south, fill=white, inner sep=1pt] at (50,33)
    {\iflangja{T2 が lock を呼ぶ}{T2 calls lock}};
  \node[lab, anchor=south, fill=white, inner sep=1pt] at (76,33)
    {\iflangja{T1 が unlock する}{T1 unlocks}};
  % (a) mutex
  \node[ttl] at (0,29) {\iflangja{(a) ミューテックス：T2 はブロック}{(a) mutex: T2 blocks}};
  \node[row] at (22,22) {CPU 1};
  \ltenseg{24}{36}{22}{white}{T1}
  \ltenseg{36}{76}{22}{ln-box}{\iflangja{T1：クリティカルセクション}{T1: critical section}}
  \ltenseg{76}{116}{22}{white}{T1}
  \node[row] at (22,15) {CPU 2};
  \ltenseg{24}{50}{15}{white}{T2}
  \ltenseg{50}{55}{15}{black!35}{}
  \ltenseg{55}{76}{15}{white}{\iflangja{T3（別のスレッド）}{T3 (another thread)}}
  \ltenseg{76}{81}{15}{black!35}{}
  \ltenseg{81}{116}{15}{ln-box}{\iflangja{T2：クリティカルセクション}{T2: critical section}}
  % (b) spinlock
  \node[ttl] at (0,6) {\iflangja{(b) スピンロック：T2 はスピン}{(b) spinlock: T2 spins}};
  \node[row] at (22,-1) {CPU 1};
  \ltenseg{24}{36}{-1}{white}{T1}
  \ltenseg{36}{76}{-1}{ln-box}{\iflangja{T1：クリティカルセクション}{T1: critical section}}
  \ltenseg{76}{116}{-1}{white}{T1}
  \node[row] at (22,-8) {CPU 2};
  \ltenseg{24}{50}{-8}{white}{T2}
  \ltenseg{50}{76}{-8}{ln-caution!20}{\iflangja{ビジーウェイト}{busy waiting}}
  \ltenseg{76}{116}{-8}{ln-box}{\iflangja{T2：クリティカルセクション}{T2: critical section}}
  % legend
  \draw[fill=black!35] (24,-18) rectangle (28,-15);
  \node[lab, anchor=west] at (29,-16.5) {\iflangja{コンテキストスイッチ}{context switch}};
  \draw[fill=ln-box] (64,-18) rectangle (68,-15);
  \node[lab, anchor=west] at (69,-16.5) {\iflangja{ロックを保持}{lock held}};
  \draw[->] (92,-16.5) -- node[lab, above] {\iflangja{時間}{time}} (116,-16.5);
\end{tikzpicture}

  \caption{別の CPU 上の T1 が保持するロックを T2 が要求する．(a)
    ミューテックスでは T2 はブロックし，CPU は別のスレッドを実行し，
    アンロックの後に T2 が戻される．(b) スピンロックでは T2 はビジー
    ウェイトし，T1 がアンロックすると同時にクリティカルセクションに入る．}
  \label{fig:l10-spin-block}
\end{figure}

したがってスピンが引き合うのは，スレッドをブロックして再開するのに要する
時間より短い時間しかロックが保持されず，かつ所有者が同時に別の CPU 上で
動作しているときだけである．単一 CPU では，スピンするスレッドが CPU を占有
している限り所有者はロックを解放できないので，スピンはタイムスライスの残り
を浪費するだけである．第5章の適応型ミューテックスは両方の振る舞いを組み
合わせ，所有者が別の CPU 上で動作しているときはスピンし，そうでなければ
ブロックする．\index{てきおうがたみゅーてっくす@適応型ミューテックス}

\begin{example}
  スレッドをブロックして後に再開するのに合計 \SI{8}{\micro\second} かかると
  する．所有者がロックを保持する時間が最大 \SI{0.5}{\micro\second} であれば，
  待っているスレッドのスピンは最大 \SI{0.5}{\micro\second} であり，ブロック
  のコストの16分の1である．これに対し所有者がロックを保持したまま
  \SI{5}{\milli\second} のディスク読み出しを行えば，スピンするスレッドは
  \SI{5}{\milli\second} の CPU 時間を，すなわちブロックのコストの
  $5000/8 = 625$ 倍を浪費する．
\end{example}

スピンロックが必要とするのは共有された1語とアトミック命令だけなので，より
複雑な機構を組み立てる基本的な構成要素となる．例えばオペレーティングシステム
のカーネルは，ミューテックスとそのブロックされたスレッドのキューの状態を
更新する短いクリティカルセクションを，スピンロックで保護する．

\section{セマフォ}
\label{sec:l10-semaphore}

Dijkstra \cite{dijkstra1965} が提案したセマフォは，オペレーティング
システムのカーネルでよく使われる同期機構である．\source{P3L4，セマフォ，
POSIX セマフォ API；Dijkstra \cite{dijkstra1965}；OSTEP ch.~31．}セマフォ
は，ミューテックスを一度に1つのスレッドから有限個のスレッドへと一般化する．

\begin{definition}[セマフォ]
  \term[せまふぉ]{セマフォ}[semaphore]とは，非負の整数で初期化される整数の
  カウントと2つの操作で表される同期機構である．\emph{ウェイト}: カウントが
  正ならそれを1減らして先へ進み，0 ならば正になるまでブロックする．
  \emph{ポスト}: カウントを1増やし，ブロックしているスレッドがあれば
  その1つを起こす．
\end{definition}

カウントは，あと何個のスレッドが通過できるかを記録する．スレッドは入るとき
ウェイトでそれを減らし，出るときポストで増やすので，初期値は同時に内部に
入れるスレッドの最大数である．2つの操作は P と V とも呼ばれ，オランダ語の
\emph{proberen}（試す）と \emph{verhogen}（増やす）に由来すると説明される
ことが多い．

セマフォはカウントの取りうる範囲で区別される．
\term[けいすうせまふぉ]{計数セマフォ}[counting semaphore]は任意の非負値を
取り，初期値と同じ数のスレッドを同時に受け入れる．カウントが 0 と 1 の値しか
取らないセマフォを\term[ばいなりせまふぉ]{バイナリセマフォ}[binary
semaphore]と呼ぶ．1 で初期化すれば一度に1つのスレッドだけを受け入れ，
ミューテックスと同じように振る舞う．

\begin{example}
  初期値 2 の計数セマフォが同一の2つのデバイスを守り，3つのスレッドがそれを
  使う．\cref{tab:l10-sem-trace}はカウントとブロックされたスレッドを追跡
  したものである．
\end{example}

\begin{table}[htbp]
  \centering
  \caption{3つのスレッドが使う初期値 2 の計数セマフォ．T3 はカウントが 0 の
    間ブロックし，T1 のポストの後に先へ進む．}
  \label{tab:l10-sem-trace}
  \small
  \begin{tabular}{@{}c>{\raggedright\arraybackslash}p{0.52\linewidth}cl@{}}
    \toprule
    手順 & 事象 & カウント & ブロック中 \\
    \midrule
    1 & T1 がウェイト: カウント $>0$，減らして進む & 1 & --- \\
    2 & T2 がウェイト: カウント $>0$，減らして進む & 0 & --- \\
    3 & T3 がウェイト: カウント $=0$，ブロック     & 0 & T3 \\
    4 & T1 がポスト: 増やして T3 を起こす          & 1 & --- \\
    5 & T3 が再検査: カウント $>0$，減らして進む   & 0 & --- \\
    6 & T2 がポスト                                & 1 & --- \\
    7 & T3 がポスト                                & 2 & --- \\
    \bottomrule
  \end{tabular}
\end{table}

\subsection{POSIX セマフォ}

\needspace{9\baselineskip}
POSIX は \texttt{<semaphore.h>} でセマフォを定義する．
\begin{lstlisting}[language=C, numbers=none]
#include <semaphore.h>

sem_t sem;
int sem_init(sem_t *sem, int pshared,
             unsigned int value);
int sem_wait(sem_t *sem);
int sem_post(sem_t *sem);
\end{lstlisting}
\texttt{sem\_init} は初期カウントを \texttt{value} に設定する．フラグ
\texttt{pshared} は，セマフォを1つのプロセスのスレッドだけが使うのか（0），
プロセス間で共有するのか（0 以外）を表す．後者の場合，\texttt{sem\_t} 自体
が，第9章の共有メモリセグメントのような，それらのプロセスが共有するメモリ
上に置かれていなければならない．\texttt{sem\_wait} と \texttt{sem\_post}
がウェイトとポストの操作である．%
\caution{セマフォに所有者はない．ウェイトしていないスレッドも含め，どの
スレッドもポストできる．}

\needspace{14\baselineskip}
\begin{example}
  あるサーバは，バックエンドのデータベースに届く要求を一度に最大4つに
  制限する．
  \begin{lstlisting}[language=C, numbers=none]
#include <semaphore.h>

#define MAX_BACKEND 4
sem_t backend_slots;

void init_slots(void) {
  sem_init(&backend_slots, 0, MAX_BACKEND);
}

void handle_request(struct request *r) {
  sem_wait(&backend_slots);   /* スロットを取る */
  forward_to_backend(r);      /* ここには最大4つ */
  sem_post(&backend_slots);   /* スロットを解放 */
}
\end{lstlisting}
  4つの要求が処理中の間に \texttt{handle\_request} を呼んだ5つ目の
  スレッドは，そのうちの1つがポストするまで \texttt{sem\_wait} で
  ブロックする．
\end{example}

\section{リーダ・ライタロック}
\label{sec:l10-rwlock}

第3章では，リーダ・ライタ方針---任意の数のリーダ，または1つのライタ---を，
ミューテックス，2つの条件変数，代理変数から組み立てた．\source{P3L4，
リーダ/ライタロック，リーダ/ライタロックの利用；OSTEP ch.~31．}この方針は
非常によく現れるので，オペレーティングシステムやスレッドライブラリはそれを
独立した機構として提供している．

\begin{definition}[リーダ・ライタロック]
  \term[りーだらいたろっく]{リーダ・ライタロック}[reader-writer lock]とは，
  2つのモードを持つロックである．\emph{読み出し}（共有）モードは任意の数の
  スレッドが同時に保持でき，\emph{書き込み}（排他）モードは，他のどの
  スレッドもどちらのモードでもロックを保持していないときに限り，1つの
  スレッドが保持できる．
\end{definition}

\needspace{12\baselineskip}
スレッドはどのモードでロックが必要かを表明し，機構が方針を強制する．Linux
カーネルでの型は \texttt{rwlock\_t} である．\margin{Linux の
\texttt{rwlock\_t} はスピンする．ブロックする変種は
\texttt{struct rw\_semaphore} である．}
\begin{lstlisting}[language=C, numbers=none]
#include <linux/spinlock.h>

rwlock_t m;

read_lock(&m);
/* クリティカルセクション: 共有状態を読む */
read_unlock(&m);

write_lock(&m);
/* クリティカルセクション: 読んで変更する */
write_unlock(&m);
\end{lstlisting}
POSIX スレッド，Java，Windows の .NET フレームワークも同じ機構を提供する．
POSIX スレッドでは次のようになる．
\begin{lstlisting}[language=C, numbers=none]
pthread_rwlock_t m;
pthread_rwlock_rdlock(&m);   /* 読み出しモードで獲得 */
pthread_rwlock_wrlock(&m);   /* 書き込みモードで獲得 */
pthread_rwlock_unlock(&m);   /* どちらのモードも解放 */
\end{lstlisting}

実装は基本的な方針では一致するが，プログラムの振る舞いに影響する細部が
異なる．
\begin{itemize}
  \item \textbf{再帰的な読み出しロック．}スレッドが読み出しロックを2回獲得
    したとき，1回の \texttt{read\_unlock} が両方の獲得を解放するのか，
    一方だけを解放するのかは，実装によって異なる．
  \item \textbf{昇格と降格．}実装によっては，リーダがロックを解放せずに
    書き込みロックへ変換したり，ライタが読み出しロックへ変換したりできる．
    これを許すかどうか，またそうした要求がすでに待っているスレッドに対して
    どの優先度を持つかは，実装によって異なる．
  \item \textbf{スケジューリングとの相互作用．}実装は，ロックが現在読み出し
    モードでしか保持されていなくても，より高い優先度かもしれないライタが
    待っている間，到着したリーダをブロックすることがある．
\end{itemize}

\begin{example}
  リーダが \SI{3}{\milli\second} ごとに到着し，各々が読み出しモードで
  \SI{5}{\milli\second} ロックを保持するとする．どのリーダも先行するリーダ
  が去る前に到着するので，ロックが空くことはない．したがって，ロックが
  読み出しモードで保持されている間つねにリーダを受け入れる実装では，待って
  いるライタは無限に飢餓状態に置かれる．ライタが待ち始めたら新しいリーダを
  ブロックする実装なら，すでに内部にいるリーダが去り次第，待ち始めてから
  最大 \SI{5}{\milli\second} でライタを受け入れる．
\end{example}

\section{モニタ}
\label{sec:l10-monitor}

ここまでの機構は，適切な場所でロックとアンロックを行う責任を，なお
プログラマに残している．\source{P3L4，モニタ；Hoare \cite{hoare1974}；
Lampson と Redell \cite{lampson1980}；Silberschatz ら ch.~6．}モニタはこの
責任を機構の側へ移す．

\begin{definition}[モニタ]
  \term[もにた]{モニタ}[monitor]とは，共有資源をカプセル化する高水準の同期
  機構である．資源と，その資源にアクセスする唯一の手段である手続き---各々を
  \term[えんとりてつづき]{エントリ手続き}[entry procedure]と呼ぶ---と，
  モニタの内部にいるスレッドが待つことのできる条件変数とを指定する．
\end{definition}

相互排除は暗黙的である（\cref{fig:l10-monitor}）．スレッドがエントリ手続き
を呼ぶと，モニタのロックが自動的に獲得され，モニタが使用中であることを
見つけたスレッドはその入口キューで待つ．手続きから戻ると，ロックは自動的に
解放される．モニタの内部で条件変数上で待つスレッドはロックを解放し，起こさ
れたときには先へ進む前にロックを再獲得する．これもまた明示的なロック操作を
伴わない．概念的には，モニタは入口でロックと検査を行い，出口でアンロックと
検査とシグナルを行う．とはいえ，古典的な設計ではシグナルそのものは明示的
であり，スレッドが待つ条件も明示的である．Hoare のモニタはウェイトとシグナル
を提供し，Mesa はこれにノーティファイとブロードキャストを加える．自動
シグナルを備えたモニタの設計はさらに進んで，スレッドがモニタを去るたびに
待機条件自体を再評価する．

\begin{example}
  容量 \texttt{N} の有界バッファをモニタとして書くと，C 風の擬似コードで
  次のようになる．
  \begin{lstlisting}[language=C, numbers=none]
monitor BoundedBuffer {
  item buf[N];
  int count = 0, in = 0, out = 0;
  condition not_full, not_empty;

  entry void put(item x) {
    while (count == N)
      wait(not_full);
    buf[in] = x;  in = (in + 1) % N;  count++;
    signal(not_empty);
  }

  entry item get(void) {
    while (count == 0)
      wait(not_empty);
    item x = buf[out];  out = (out + 1) % N;
    count--;
    signal(not_full);
    return x;
  }
}
\end{lstlisting}
  コードにロック操作は現れないので，忘れることも，誤ったロックに適用する
  こともできない．
\end{example}

\begin{figure}[htbp]
  \centering
  % Structure of a monitor: entry queue, implicit lock, entry procedures,
% shared data and condition-variable queues.
% Usage: % Structure of a monitor: entry queue, implicit lock, entry procedures,
% shared data and condition-variable queues.
% Usage: % Structure of a monitor: entry queue, implicit lock, entry procedures,
% shared data and condition-variable queues.
% Usage: \input{figures/l10-monitor}   (width ~116 mm)
\begin{tikzpicture}[x=1mm, y=1mm, font=\footnotesize\sffamily,
  box/.style={draw, rounded corners=2pt, fill=white, align=center, inner sep=2pt},
  q/.style={draw, rounded corners=2pt, fill=ln-box, minimum width=22mm,
            minimum height=7mm, inner sep=1pt},
  th/.style={draw, circle, fill=white, minimum size=4.2mm, inner sep=0pt,
             font=\tiny\sffamily},
  lab/.style={font=\scriptsize\sffamily, align=center},
  >={Stealth[length=1.8mm]}]
  % monitor
  \draw[thick, rounded corners=4pt, fill=ln-box] (32,-17) rectangle (84,19);
  \node[font=\footnotesize\sffamily\bfseries, anchor=north west] at (33,18.5)
    {\iflangja{モニタ}{monitor}};
  \node[box, minimum width=44mm, minimum height=7mm] (data) at (58,6)
    {\iflangja{共有データ}{shared data}\quad\texttt{buf}, \texttt{count}};
  \node[box, minimum width=19mm, minimum height=6mm, font=\footnotesize\ttfamily] (put) at (46.5,-4) {put()};
  \node[box, minimum width=19mm, minimum height=6mm, font=\footnotesize\ttfamily] (get) at (69.5,-4) {get()};
  \node[lab, anchor=north] at (58,-8.5) {\iflangja{エントリ手続き}{entry procedures}\\
    \iflangja{（内部で動くのは一度に1スレッド）}{(one thread inside at a time)}};
  \draw[ln-muted] (put.north) -- (data.south -| put.north);
  \draw[ln-muted] (get.north) -- (data.south -| get.north);
  % entry queue
  \node[q] (eq) at (13,0) {};
  \foreach \i/\t in {0/T5,1/T4} { \node[th] at (7.5+6*\i,0) {\t}; }
  \node[lab, anchor=south] at (13,3.5) {\iflangja{入口キュー}{entry queue}};
  \draw[->] (eq.east) -- node[lab, above, pos=0.45] {\iflangja{ロック}{lock}} (32,0);
  % condition queues
  \node[q] (nf) at (105,9) {};
  \node[th] at (99.5,9) {T2};
  \node[lab, anchor=south] at (105,12.5) {\texttt{not\_full}};
  \node[q] (ne) at (105,-9) {};
  \node[th] at (99.5,-9) {T3};
  \node[lab, anchor=south] at (105,-5.5) {\texttt{not\_empty}};
  \draw[->] (84,9) -- (nf.west);
  \draw[->] (84,-9) -- (ne.west);
  \node[lab, anchor=south] at (89,11.5) {\texttt{wait}};
  \node[lab] at (90,0) {\iflangja{ロック\\を解放}{releases\\lock}};
  % signal path back to the entry queue
  \draw[->] (ne.south) -- ++(0,-12) -| node[lab, pos=0.25, below]
    {\texttt{signal} / \texttt{broadcast}:
     \iflangja{起こされたスレッドはロックを再獲得する}{a woken thread re-acquires the lock}} (eq.south);
  % exit
  \draw[->] (58,19) -- ++(0,6) node[lab, anchor=south] {\iflangja{return（ロックを解放）}{return (releases lock)}};
\end{tikzpicture}
   (width ~116 mm)
\begin{tikzpicture}[x=1mm, y=1mm, font=\footnotesize\sffamily,
  box/.style={draw, rounded corners=2pt, fill=white, align=center, inner sep=2pt},
  q/.style={draw, rounded corners=2pt, fill=ln-box, minimum width=22mm,
            minimum height=7mm, inner sep=1pt},
  th/.style={draw, circle, fill=white, minimum size=4.2mm, inner sep=0pt,
             font=\tiny\sffamily},
  lab/.style={font=\scriptsize\sffamily, align=center},
  >={Stealth[length=1.8mm]}]
  % monitor
  \draw[thick, rounded corners=4pt, fill=ln-box] (32,-17) rectangle (84,19);
  \node[font=\footnotesize\sffamily\bfseries, anchor=north west] at (33,18.5)
    {\iflangja{モニタ}{monitor}};
  \node[box, minimum width=44mm, minimum height=7mm] (data) at (58,6)
    {\iflangja{共有データ}{shared data}\quad\texttt{buf}, \texttt{count}};
  \node[box, minimum width=19mm, minimum height=6mm, font=\footnotesize\ttfamily] (put) at (46.5,-4) {put()};
  \node[box, minimum width=19mm, minimum height=6mm, font=\footnotesize\ttfamily] (get) at (69.5,-4) {get()};
  \node[lab, anchor=north] at (58,-8.5) {\iflangja{エントリ手続き}{entry procedures}\\
    \iflangja{（内部で動くのは一度に1スレッド）}{(one thread inside at a time)}};
  \draw[ln-muted] (put.north) -- (data.south -| put.north);
  \draw[ln-muted] (get.north) -- (data.south -| get.north);
  % entry queue
  \node[q] (eq) at (13,0) {};
  \foreach \i/\t in {0/T5,1/T4} { \node[th] at (7.5+6*\i,0) {\t}; }
  \node[lab, anchor=south] at (13,3.5) {\iflangja{入口キュー}{entry queue}};
  \draw[->] (eq.east) -- node[lab, above, pos=0.45] {\iflangja{ロック}{lock}} (32,0);
  % condition queues
  \node[q] (nf) at (105,9) {};
  \node[th] at (99.5,9) {T2};
  \node[lab, anchor=south] at (105,12.5) {\texttt{not\_full}};
  \node[q] (ne) at (105,-9) {};
  \node[th] at (99.5,-9) {T3};
  \node[lab, anchor=south] at (105,-5.5) {\texttt{not\_empty}};
  \draw[->] (84,9) -- (nf.west);
  \draw[->] (84,-9) -- (ne.west);
  \node[lab, anchor=south] at (89,11.5) {\texttt{wait}};
  \node[lab] at (90,0) {\iflangja{ロック\\を解放}{releases\\lock}};
  % signal path back to the entry queue
  \draw[->] (ne.south) -- ++(0,-12) -| node[lab, pos=0.25, below]
    {\texttt{signal} / \texttt{broadcast}:
     \iflangja{起こされたスレッドはロックを再獲得する}{a woken thread re-acquires the lock}} (eq.south);
  % exit
  \draw[->] (58,19) -- ++(0,6) node[lab, anchor=south] {\iflangja{return（ロックを解放）}{return (releases lock)}};
\end{tikzpicture}
   (width ~116 mm)
\begin{tikzpicture}[x=1mm, y=1mm, font=\footnotesize\sffamily,
  box/.style={draw, rounded corners=2pt, fill=white, align=center, inner sep=2pt},
  q/.style={draw, rounded corners=2pt, fill=ln-box, minimum width=22mm,
            minimum height=7mm, inner sep=1pt},
  th/.style={draw, circle, fill=white, minimum size=4.2mm, inner sep=0pt,
             font=\tiny\sffamily},
  lab/.style={font=\scriptsize\sffamily, align=center},
  >={Stealth[length=1.8mm]}]
  % monitor
  \draw[thick, rounded corners=4pt, fill=ln-box] (32,-17) rectangle (84,19);
  \node[font=\footnotesize\sffamily\bfseries, anchor=north west] at (33,18.5)
    {\iflangja{モニタ}{monitor}};
  \node[box, minimum width=44mm, minimum height=7mm] (data) at (58,6)
    {\iflangja{共有データ}{shared data}\quad\texttt{buf}, \texttt{count}};
  \node[box, minimum width=19mm, minimum height=6mm, font=\footnotesize\ttfamily] (put) at (46.5,-4) {put()};
  \node[box, minimum width=19mm, minimum height=6mm, font=\footnotesize\ttfamily] (get) at (69.5,-4) {get()};
  \node[lab, anchor=north] at (58,-8.5) {\iflangja{エントリ手続き}{entry procedures}\\
    \iflangja{（内部で動くのは一度に1スレッド）}{(one thread inside at a time)}};
  \draw[ln-muted] (put.north) -- (data.south -| put.north);
  \draw[ln-muted] (get.north) -- (data.south -| get.north);
  % entry queue
  \node[q] (eq) at (13,0) {};
  \foreach \i/\t in {0/T5,1/T4} { \node[th] at (7.5+6*\i,0) {\t}; }
  \node[lab, anchor=south] at (13,3.5) {\iflangja{入口キュー}{entry queue}};
  \draw[->] (eq.east) -- node[lab, above, pos=0.45] {\iflangja{ロック}{lock}} (32,0);
  % condition queues
  \node[q] (nf) at (105,9) {};
  \node[th] at (99.5,9) {T2};
  \node[lab, anchor=south] at (105,12.5) {\texttt{not\_full}};
  \node[q] (ne) at (105,-9) {};
  \node[th] at (99.5,-9) {T3};
  \node[lab, anchor=south] at (105,-5.5) {\texttt{not\_empty}};
  \draw[->] (84,9) -- (nf.west);
  \draw[->] (84,-9) -- (ne.west);
  \node[lab, anchor=south] at (89,11.5) {\texttt{wait}};
  \node[lab] at (90,0) {\iflangja{ロック\\を解放}{releases\\lock}};
  % signal path back to the entry queue
  \draw[->] (ne.south) -- ++(0,-12) -| node[lab, pos=0.25, below]
    {\texttt{signal} / \texttt{broadcast}:
     \iflangja{起こされたスレッドはロックを再獲得する}{a woken thread re-acquires the lock}} (eq.south);
  % exit
  \draw[->] (58,19) -- ++(0,6) node[lab, anchor=south] {\iflangja{return（ロックを解放）}{return (releases lock)}};
\end{tikzpicture}

  \caption{モニタ．スレッドはエントリ手続きを通じてのみ内部に入り，
    エントリ手続きはモニタのロックを暗黙的に獲得する．条件変数上で待つ
    スレッドはロックを解放し，ここで示した Mesa のセマンティクスでは，
    起こされた後にロックを再獲得しなければならない．}
  \label{fig:l10-monitor}
\end{figure}

モニタは Xerox PARC で開発された Mesa 言語の一部であった
\cite{lampson1980}．\margin{Mesa と Java では，起こされたスレッドは
\texttt{while} ループで条件を再検査する．}今日では Java がモニタを
サポートしている．すべてのオブジェクトが内部にロックを持ち，
\texttt{synchronized} と宣言されたメソッドは，そのロックとアンロックを
言語ランタイムが行うエントリ手続きである．ただし通知は \texttt{notify} と
\texttt{notifyAll} の明示的な呼び出しのままである．

\emph{モニタ}という語はプログラミングのスタイルも指す．第3章のように，
ミューテックスと条件変数から組み立てた明示的な入口コードと出口コードで
すべてのクリティカルセクションを構成することは，言語の支援なしにモニタの
規律を適用することである．

\section{その他の同期機構}
\label{sec:l10-other}

同期方針をより述べやすくすること，あるいは待つこと自体を避けることを目指す
機構が，ほかにもいくつかある．\source{P3L4，さらなる同期機構．}
\begin{description}
  \item[シリアライザ] \term[しりあらいざ]{シリアライザ}[serializer]
    \cite{hewitt1979}はモニタと同様に資源を守るが，優先度を定義しやすくし，
    明示的なシグナルを隠す．スレッドは，自分が先へ進める条件とともにキュー
    で待つ．条件が成立し，かつそのスレッドが自分のキューの先頭にあるとき，
    シリアライザが自動的にそれを再開する．したがってプログラムには条件変数
    もシグナル呼び出しも現れず，キューの順序が優先度を表す．
  \item[パス式] \term[ぱすしき]{パス式}[path expression]\cite{campbell1974}
    は，資源に対する操作の名前を用いた正規表現として，操作が実行されうる
    順序を述べる．ランタイムシステムは，その式が許す順序でのみ操作を交錯
    させる．「多数の読み出しまたは1つの書き込み」という方針は
    \texttt{path \{read\}, write end} と書く．波括弧は任意個の読み出しの
    同時実行を許し，コンマは選択肢の1つを選び，パスは無限に繰り返される．
  \item[バリア] \term[ばりあ]{バリア}[barrier]はセマフォの逆である．
    セマフォは最大 $n$ 個のスレッドを通し，その後に来たスレッドをブロック
    する．バリアは $n$ 個のスレッドが到着するまでスレッドをブロックし，
    その後すべてをまとめて進ませる（\cref{fig:l10-barrier}）．POSIX
    スレッドは \texttt{pthread\_barrier\_t} と
    \texttt{pthread\_barrier\_wait} を提供する．
  \item[ランデブーポイント]
    \term[らんでぶーぽいんと]{ランデブーポイント}[rendezvous point]とは，
    2つ以上の特定のスレッドの実行における点であり，各スレッドは他の
    スレッドが対応する点に達するまでそこで待つ．その後，典型的には
    データを交換してから，スレッドは実行を続ける．
\end{description}

\begin{figure}[htbp]
  \centering
  % A barrier for n = 3 threads: threads block at the barrier until the
% third one arrives, then all continue.
% Usage: % A barrier for n = 3 threads: threads block at the barrier until the
% third one arrives, then all continue.
% Usage: % A barrier for n = 3 threads: threads block at the barrier until the
% third one arrives, then all continue.
% Usage: \input{figures/l10-barrier}   (width ~116 mm)
\begin{tikzpicture}[x=1mm, y=1mm, font=\footnotesize\sffamily,
  lab/.style={font=\scriptsize\sffamily, align=center},
  run/.style={line width=1.6pt, ln-accent},
  blk/.style={line width=4pt, ln-caution!25},
  >={Stealth[length=1.8mm]}]
  \foreach \y/\t/\a in {10/T1/40, 3/T2/62, -4/T3/82} {
    \node[lab, anchor=east] at (18,\y) {\t};
    \draw[run] (20,\y) -- (\a,\y);
    \ifnum\a<82 \draw[blk] (\a,\y) -- (82,\y); \fi
    \draw[run, ->] (82,\y) -- (112,\y);
    \fill[ln-accent] (\a,\y) circle (0.9);
  }
  \draw[line width=1.2pt] (82,15) -- (82,-9);
  \node[lab, anchor=south] at (82,15)
    {\iflangja{バリア（$n=3$）}{barrier ($n=3$)}};
  \node[lab, anchor=south] at (40,12)
    {\iflangja{T1 が到着}{T1 arrives}};
  \node[lab, anchor=west] at (83,-7.5)
    {\iflangja{T3 の到着で全員が進む}{T3 arrives: all proceed}};
  % legend and time axis
  \draw[run] (20,-15) -- (28,-15);
  \node[lab, anchor=west] at (29,-15) {\iflangja{実行中}{running}};
  \draw[blk] (46,-15) -- (54,-15);
  \node[lab, anchor=west] at (55,-15) {\iflangja{バリアで待機}{blocked at the barrier}};
  \draw[->] (92,-15) -- node[lab, above] {\iflangja{時間}{time}} (116,-15);
\end{tikzpicture}
   (width ~116 mm)
\begin{tikzpicture}[x=1mm, y=1mm, font=\footnotesize\sffamily,
  lab/.style={font=\scriptsize\sffamily, align=center},
  run/.style={line width=1.6pt, ln-accent},
  blk/.style={line width=4pt, ln-caution!25},
  >={Stealth[length=1.8mm]}]
  \foreach \y/\t/\a in {10/T1/40, 3/T2/62, -4/T3/82} {
    \node[lab, anchor=east] at (18,\y) {\t};
    \draw[run] (20,\y) -- (\a,\y);
    \ifnum\a<82 \draw[blk] (\a,\y) -- (82,\y); \fi
    \draw[run, ->] (82,\y) -- (112,\y);
    \fill[ln-accent] (\a,\y) circle (0.9);
  }
  \draw[line width=1.2pt] (82,15) -- (82,-9);
  \node[lab, anchor=south] at (82,15)
    {\iflangja{バリア（$n=3$）}{barrier ($n=3$)}};
  \node[lab, anchor=south] at (40,12)
    {\iflangja{T1 が到着}{T1 arrives}};
  \node[lab, anchor=west] at (83,-7.5)
    {\iflangja{T3 の到着で全員が進む}{T3 arrives: all proceed}};
  % legend and time axis
  \draw[run] (20,-15) -- (28,-15);
  \node[lab, anchor=west] at (29,-15) {\iflangja{実行中}{running}};
  \draw[blk] (46,-15) -- (54,-15);
  \node[lab, anchor=west] at (55,-15) {\iflangja{バリアで待機}{blocked at the barrier}};
  \draw[->] (92,-15) -- node[lab, above] {\iflangja{時間}{time}} (116,-15);
\end{tikzpicture}
   (width ~116 mm)
\begin{tikzpicture}[x=1mm, y=1mm, font=\footnotesize\sffamily,
  lab/.style={font=\scriptsize\sffamily, align=center},
  run/.style={line width=1.6pt, ln-accent},
  blk/.style={line width=4pt, ln-caution!25},
  >={Stealth[length=1.8mm]}]
  \foreach \y/\t/\a in {10/T1/40, 3/T2/62, -4/T3/82} {
    \node[lab, anchor=east] at (18,\y) {\t};
    \draw[run] (20,\y) -- (\a,\y);
    \ifnum\a<82 \draw[blk] (\a,\y) -- (82,\y); \fi
    \draw[run, ->] (82,\y) -- (112,\y);
    \fill[ln-accent] (\a,\y) circle (0.9);
  }
  \draw[line width=1.2pt] (82,15) -- (82,-9);
  \node[lab, anchor=south] at (82,15)
    {\iflangja{バリア（$n=3$）}{barrier ($n=3$)}};
  \node[lab, anchor=south] at (40,12)
    {\iflangja{T1 が到着}{T1 arrives}};
  \node[lab, anchor=west] at (83,-7.5)
    {\iflangja{T3 の到着で全員が進む}{T3 arrives: all proceed}};
  % legend and time axis
  \draw[run] (20,-15) -- (28,-15);
  \node[lab, anchor=west] at (29,-15) {\iflangja{実行中}{running}};
  \draw[blk] (46,-15) -- (54,-15);
  \node[lab, anchor=west] at (55,-15) {\iflangja{バリアで待機}{blocked at the barrier}};
  \draw[->] (92,-15) -- node[lab, above] {\iflangja{時間}{time}} (116,-15);
\end{tikzpicture}

  \caption{3つのスレッドのためのバリア．T1 と T2 はバリアに達したところで
    ブロックし，T3 の到着が3つすべてを解放する．}
  \label{fig:l10-barrier}
\end{figure}

\subsection{ウェイトフリー同期と RCU}

ここまでの機構はいずれもスレッドを待たせる．楽観的な代替手段は，並行する
操作が衝突しないことに賭けてスレッドをロックなしに進ませ，変更されている
間も並行するリーダが安全であるようにデータ構造を設計する．

\begin{definition}[ウェイトフリー同期]
  他のスレッドがどれだけ速くても遅くても，どのスレッドも各操作を自分自身の
  有限回のステップで完了する同期を
  \term[うぇいとふりーどうき]{ウェイトフリー同期}[wait-free
  synchronization] \cite{herlihy1991}と呼ぶ．どのスレッドも他のスレッドを
  待つことがない．
\end{definition}

現在は Linux カーネルの一部である \term{RCU}[read-copy-update] 機構
\cite{mckenney1998} は，この考えを，書かれるよりはるかに多く読まれるデータ
に適用する．リーダはロックをまったく取らない．ライタは現在の版を複製し，
複製を変更し，リーダがデータを見つけるために辿るポインタをアトミックに
差し替えて公開する．そのためリーダは，古い版か新しい版のどちらかを見るので
あって，中途半端に変更された版を見ることはない．古い版は，なおそれへの参照
を保持しているかもしれないリーダが終えた時点で解放される．したがってリーダ
は決して待たず，ライタどうしは通常のロックで協調する．

\cref{tab:l10-constructs}はこの章の機構をミューテックスと比較したもので
ある．スレッドを待たせるかどうかにかかわらず，そのいずれもがメモリ位置を
アトミックに更新すること---ロック語の設定，カウントの増減，ポインタの
差し替え---に依存しており，これにはハードウェアの支援が必要である．

\begin{table}[htbp]
  \centering
  \caption{同期機構の比較．}
  \label{tab:l10-constructs}
  \small
  \begin{tabular}{@{}lp{32mm}p{30mm}l@{}}
    \toprule
    機構 & 同時に受け入れる数 & 待つスレッド & ロック \\
    \midrule
    ミューテックス       & 1つ                      & ブロック   & 明示的 \\
    スピンロック         & 1つ                      & スピン     & 明示的 \\
    計数セマフォ         & 初期カウントまで         & ブロック   & 明示的 \\
    リーダ・ライタロック & 任意個のリーダまたは1つのライタ
                                                    & ブロックまたはスピン
                                                                 & 明示的 \\
    モニタ               & 1つ                      & ブロック   & 暗黙的 \\
    バリア               & $n$ 個到着したら $n$ 個  & ブロック   & 明示的 \\
    RCU                  & 任意個のリーダと1つのライタ
                                                    & リーダは待たない
                                                                 & ライタのみ \\
    \bottomrule
  \end{tabular}
\end{table}

\begin{keypoint}
  スピンロックはコンテキストスイッチと引き換えに CPU 時間を費やし，セマフォは
  有限個のスレッドを受け入れ，リーダ・ライタロックはリーダ・ライタ方針を
  1つの機構にまとめ，モニタはロックを暗黙的にする．シリアライザ，パス式，
  バリア，ランデブーポイントはさらに別の方針を表現し，RCU のような
  ウェイトフリーの方式は待つこと自体を避ける．
\end{keypoint}

\section{同期のためのハードウェア支援}
\label{sec:l10-hw}

まず，教科書がロックの最初の試みとして挙げる形で，通常の命令だけでスピン
ロックを実装してみる．\source{P3L4，同期の構成要素: スピンロック，
ハードウェア支援の必要性，アトミック命令；OSTEP ch.~28；
Silberschatz ら ch.~6．}
\needspace{10\baselineskip}
\begin{lstlisting}[language=C, numbers=none]
#define FREE 0
#define BUSY 1

void naive_lock(int *lock) {
  while (*lock == BUSY)
    ;                /* スピン */
  *lock = BUSY;      /* 検査とは別の命令 */
}
\end{lstlisting}
検査と更新は別々の命令であり，その間に別のスレッドが同じ空きロックを観測
しうる．\cref{tab:l10-naive-race}は，異なる CPU 上の2つのスレッドがともに
クリティカルセクションに入る様子を示している．2つの手順を別のロックで保護
しても，問題をそのロックへ移すだけである．この形のロックが正しいのは，
それを検査することと設定することが1つの不可分な手順として起こる場合だけで
あり，実用的な実装はその手順をハードウェアから得る．\margin{Peterson の
アルゴリズムのようなソフトウェアのアルゴリズムは，通常の読み書きだけで
相互排除を達成する．}

\begin{table}[htbp]
  \centering
  \caption{異なる CPU 上で \texttt{naive\_lock} を実行する2つのスレッド．
    どちらかが書き込む前に両方がロックを読むので，両方がクリティカル
    セクションに入る．}
  \label{tab:l10-naive-race}
  \small
  \begin{tabular}{@{}c>{\raggedright\arraybackslash}p{0.34\linewidth}>{\raggedright\arraybackslash}p{0.34\linewidth}l@{}}
    \toprule
    手順 & CPU 1 上の T1 & CPU 2 上の T2 & \texttt{*lock} \\
    \midrule
    1 & \texttt{FREE} を読み，ループを抜ける &                                      & \texttt{FREE} \\
    2 &                                      & \texttt{FREE} を読み，ループを抜ける & \texttt{FREE} \\
    3 & \texttt{BUSY} を書き込み，中に入る   &                                      & \texttt{BUSY} \\
    4 &                                      & \texttt{BUSY} を書き込み，中に入る   & \texttt{BUSY} \\
    \bottomrule
  \end{tabular}
\end{table}

\begin{definition}[アトミック命令]
  \term[あとみっくめいれい]{アトミック命令}[atomic instruction]とは，複数の
  操作，典型的にはメモリ位置の読み出しと書き込みを，1つの不可分な手順として
  実行するプロセッサ命令である．ハードウェアは，アトミック性（複数の操作が
  まとめて効力を持つこと），相互排除（同じメモリ位置に対する他のアトミック
  命令がその間に割り込まないこと），および同じ位置に対する複数の並行実行の
  うち1つを除くすべてが待たされることを保証する．
\end{definition}

したがってアトミック命令は，ハードウェアによって保護されたクリティカル
セクションである．よくある例は次のものである．
\begin{description}
  \item[\texttt{test\_and\_set}]
    \term[てすとあんどせっと]{テストアンドセット}[test-and-set]命令は，
    ロック語の現在の値を返し，その語を \texttt{BUSY} に設定する．
  \item[\texttt{read\_and\_increment}] メモリ位置の値を返し，それを1増やす．
  \item[\texttt{compare\_and\_swap}]
    \term[こんぺああんどすわっぷ]{コンペアアンドスワップ}[compare-and-swap]%
    命令は，メモリ位置を期待値と比較し，等しい場合に限り新しい値を書き込む．
    書き込みが行われたかどうかは，この命令が報告する．
\end{description}

\needspace{14\baselineskip}
プロセッサがアトミックに実行する C 関数として書けば，最初のものと最後の
ものは次のようになる．
\begin{lstlisting}[language=C, numbers=none]
/* 各関数は1つのアトミックな手順として実行される */
int test_and_set(int *lock) {
  int old = *lock;
  *lock = BUSY;
  return old;
}

int compare_and_swap(int *p, int expected,
                     int new_val) {
  if (*p != expected)
    return 0;            /* 値が変わっている */
  *p = new_val;
  return 1;
}
\end{lstlisting}

\needspace{9\baselineskip}
\texttt{test\_and\_set} を使えばスピンロックは正しくなる．
\begin{lstlisting}[language=C, numbers=none]
void spinlock_lock(int *lock) {
  while (test_and_set(lock) == BUSY)
    ;   /* スピン */
}

void spinlock_unlock(int *lock) {
  *lock = FREE;
}
\end{lstlisting}
ロックが空いていれば \texttt{test\_and\_set} は \texttt{FREE} を返し，
すでにロックをビジーと印を付けているので，スレッドは所有者としてループを
抜ける．ロックがビジーであれば，\texttt{BUSY} に \texttt{BUSY} を上書き
して---これは何も変えない---\texttt{BUSY} を返すので，スレッドは再び
試みる．空いているロックに対して同時に \texttt{test\_and\_set} を実行
した複数のスレッドのうち，\texttt{FREE} を得るのはちょうど1つである．

\begin{example}
  コンペアアンドスワップは，ロックなしに共有データを更新する．カウンタは，
  スワップが成功するまで読み直すことで1増やされる．
  \begin{lstlisting}[language=C, numbers=none]
void increment(int *counter) {
  int old;
  do {
    old = *counter;
  } while (!compare_and_swap(counter, old, old + 1));
}
\end{lstlisting}
  スワップが失敗するのは，読み出しとスワップの間に別のスレッドが
  カウンタ\caution{このループはブロックしないが，スレッドが無限に再試行
  しなければならないこともあるので，ウェイトフリーではない．}を変更した
  ときに限られる．そのときループは読み直して再試行するので，増分が失われる
  ことはない．
\end{example}

\subsection{移植性}

アトミック命令はプロセッサアーキテクチャに固有である．どの命令が存在し，
それがどれだけ効率的かはプラットフォームによって異なる．例えば x86 では，
メモリオペランドを取る \texttt{xchg} がテストアンドセットとして，
\texttt{lock cmpxchg} がコンペアアンドスワップとして使える．したがって，
カーネルやスレッドライブラリのように同期機構を実装するソフトウェアは，
動作するすべてのアーキテクチャへ移植しなければならない．さらにプラット
フォームごとに特化もされる．同じ機構の複数の実装のうち，そのプラット
フォームが提供する最も効率的な命令の上に組み立てられたものを，条件
コンパイルで選ぶ．

\section{共有メモリ型マルチプロセッサとキャッシュ}
\label{sec:l10-smp}

同期機構が最も重要になるのは，スレッドが実際に同時に動作する複数 CPU の
マシン上である．\source{P3L4，共有メモリ型マルチプロセッサ，キャッシュ
一貫性，キャッシュ一貫性とアトミック命令；Silberschatz ら ch.~1．}共有
メモリ型マルチプロセッサ（第7章）%
\index{きょうゆうめもりがたまるちぷろせっさ@共有メモリ型マルチプロセッサ}%
は，複数の CPU と，そのすべてがアクセスできるメモリを持つ．こうしたシステム
は対称型マルチプロセッサ（SMP）とも呼ばれる．\cref{fig:l10-smp}は2つの構成
を示す．古いシステムに典型的なバス型の構成では，CPU と単一のメモリ部品が
1本の共有バスに接続され，そのバスは一度に1つのメモリアクセスしか運べない．
今日一般的なインターコネクト型の構成では，CPU と，場合によっては複数の
メモリ部品が，複数のメモリアクセスを同時に進行させられるインターコネクトを
通じて接続される．

\begin{figure}[htbp]
  \centering
  % Shared memory multiprocessor configurations: bus-based and
% interconnect-based, each CPU with its own cache.
% Usage: % Shared memory multiprocessor configurations: bus-based and
% interconnect-based, each CPU with its own cache.
% Usage: % Shared memory multiprocessor configurations: bus-based and
% interconnect-based, each CPU with its own cache.
% Usage: \input{figures/l10-smp}   (width ~116 mm)
\begin{tikzpicture}[x=1mm, y=1mm, font=\footnotesize\sffamily,
  cpu/.style={draw, fill=white, minimum width=13mm, minimum height=6mm, inner sep=1pt},
  cache/.style={draw, fill=ln-box, minimum width=13mm, minimum height=4.5mm,
                inner sep=1pt, font=\scriptsize\sffamily},
  mem/.style={draw, fill=ln-box, minimum height=7mm, inner sep=1pt},
  lab/.style={font=\scriptsize\sffamily, align=center}]
  % (a) bus-based
  \foreach \x/\n in {8/1, 27/2, 46/3} {
    \node[cpu] (a\n) at (\x,0) {CPU \n};
    \node[cache, anchor=north] (ac\n) at (a\n.south) {\iflangja{キャッシュ}{cache}};
    \draw (ac\n.south) -- (\x,-13);
  }
  \draw[line width=1.6pt] (0,-13) -- (54,-13);
  \node[lab, anchor=south west] at (48,-12.6) {\iflangja{バス}{bus}};
  \node[mem, minimum width=30mm] (am) at (27,-22) {\iflangja{メモリ}{memory}};
  \draw (am.north) -- (27,-13);
  \node[lab] at (27,-31) {\iflangja{(a) バス型}{(a) bus-based}};
  % (b) interconnect-based
  \foreach \x/\n in {70/1, 89/2, 108/3} {
    \node[cpu] (b\n) at (\x,0) {CPU \n};
    \node[cache, anchor=north] (bc\n) at (b\n.south) {\iflangja{キャッシュ}{cache}};
    \draw (bc\n.south) -- (\x,-10.5);
  }
  \node[draw, fill=white, minimum width=50mm, minimum height=5mm, inner sep=1pt]
    (ic) at (89,-13) {\iflangja{インターコネクト}{interconnect}};
  \node[mem, minimum width=22mm] (bm1) at (76,-22) {\iflangja{メモリ}{memory}};
  \node[mem, minimum width=22mm] (bm2) at (102,-22) {\iflangja{メモリ}{memory}};
  \draw (bm1.north) -- (76,-15.5);
  \draw (bm2.north) -- (102,-15.5);
  \node[lab] at (89,-31) {\iflangja{(b) インターコネクト型}{(b) interconnect-based}};
\end{tikzpicture}
   (width ~116 mm)
\begin{tikzpicture}[x=1mm, y=1mm, font=\footnotesize\sffamily,
  cpu/.style={draw, fill=white, minimum width=13mm, minimum height=6mm, inner sep=1pt},
  cache/.style={draw, fill=ln-box, minimum width=13mm, minimum height=4.5mm,
                inner sep=1pt, font=\scriptsize\sffamily},
  mem/.style={draw, fill=ln-box, minimum height=7mm, inner sep=1pt},
  lab/.style={font=\scriptsize\sffamily, align=center}]
  % (a) bus-based
  \foreach \x/\n in {8/1, 27/2, 46/3} {
    \node[cpu] (a\n) at (\x,0) {CPU \n};
    \node[cache, anchor=north] (ac\n) at (a\n.south) {\iflangja{キャッシュ}{cache}};
    \draw (ac\n.south) -- (\x,-13);
  }
  \draw[line width=1.6pt] (0,-13) -- (54,-13);
  \node[lab, anchor=south west] at (48,-12.6) {\iflangja{バス}{bus}};
  \node[mem, minimum width=30mm] (am) at (27,-22) {\iflangja{メモリ}{memory}};
  \draw (am.north) -- (27,-13);
  \node[lab] at (27,-31) {\iflangja{(a) バス型}{(a) bus-based}};
  % (b) interconnect-based
  \foreach \x/\n in {70/1, 89/2, 108/3} {
    \node[cpu] (b\n) at (\x,0) {CPU \n};
    \node[cache, anchor=north] (bc\n) at (b\n.south) {\iflangja{キャッシュ}{cache}};
    \draw (bc\n.south) -- (\x,-10.5);
  }
  \node[draw, fill=white, minimum width=50mm, minimum height=5mm, inner sep=1pt]
    (ic) at (89,-13) {\iflangja{インターコネクト}{interconnect}};
  \node[mem, minimum width=22mm] (bm1) at (76,-22) {\iflangja{メモリ}{memory}};
  \node[mem, minimum width=22mm] (bm2) at (102,-22) {\iflangja{メモリ}{memory}};
  \draw (bm1.north) -- (76,-15.5);
  \draw (bm2.north) -- (102,-15.5);
  \node[lab] at (89,-31) {\iflangja{(b) インターコネクト型}{(b) interconnect-based}};
\end{tikzpicture}
   (width ~116 mm)
\begin{tikzpicture}[x=1mm, y=1mm, font=\footnotesize\sffamily,
  cpu/.style={draw, fill=white, minimum width=13mm, minimum height=6mm, inner sep=1pt},
  cache/.style={draw, fill=ln-box, minimum width=13mm, minimum height=4.5mm,
                inner sep=1pt, font=\scriptsize\sffamily},
  mem/.style={draw, fill=ln-box, minimum height=7mm, inner sep=1pt},
  lab/.style={font=\scriptsize\sffamily, align=center}]
  % (a) bus-based
  \foreach \x/\n in {8/1, 27/2, 46/3} {
    \node[cpu] (a\n) at (\x,0) {CPU \n};
    \node[cache, anchor=north] (ac\n) at (a\n.south) {\iflangja{キャッシュ}{cache}};
    \draw (ac\n.south) -- (\x,-13);
  }
  \draw[line width=1.6pt] (0,-13) -- (54,-13);
  \node[lab, anchor=south west] at (48,-12.6) {\iflangja{バス}{bus}};
  \node[mem, minimum width=30mm] (am) at (27,-22) {\iflangja{メモリ}{memory}};
  \draw (am.north) -- (27,-13);
  \node[lab] at (27,-31) {\iflangja{(a) バス型}{(a) bus-based}};
  % (b) interconnect-based
  \foreach \x/\n in {70/1, 89/2, 108/3} {
    \node[cpu] (b\n) at (\x,0) {CPU \n};
    \node[cache, anchor=north] (bc\n) at (b\n.south) {\iflangja{キャッシュ}{cache}};
    \draw (bc\n.south) -- (\x,-10.5);
  }
  \node[draw, fill=white, minimum width=50mm, minimum height=5mm, inner sep=1pt]
    (ic) at (89,-13) {\iflangja{インターコネクト}{interconnect}};
  \node[mem, minimum width=22mm] (bm1) at (76,-22) {\iflangja{メモリ}{memory}};
  \node[mem, minimum width=22mm] (bm2) at (102,-22) {\iflangja{メモリ}{memory}};
  \draw (bm1.north) -- (76,-15.5);
  \draw (bm2.north) -- (102,-15.5);
  \node[lab] at (89,-31) {\iflangja{(b) インターコネクト型}{(b) interconnect-based}};
\end{tikzpicture}

  \caption{共有メモリ型マルチプロセッサの構成．どちらでも，各 CPU が自分の
    キャッシュを持ち，すべての CPU がすべてのメモリにアクセスできる．}
  \label{fig:l10-smp}
\end{figure}

どの CPU も自分のキャッシュを持つ．メモリはキャッシュよりも到達に時間が
かかる．物理的により遠いからでもあり，そこへ至るバスやインターコネクトを
CPU どうしが奪い合うからでもある．キャッシュはこのレイテンシを隠す．

しかし複数の CPU が同じデータを扱うと，1つのメモリ位置の複製が複数の
キャッシュに同時に存在しうる．ある CPU がその位置に書き込むと，他の CPU が
持つ複製は古くなり，そうした複製を読んだスレッドは古い値を見てしまう
（\cref{fig:l10-coherence}）．メモリ位置の複製が複数の CPU にキャッシュ
されていても，どの CPU もその位置の最新の値を観測するという性質を，
\term[きゃっしゅいっかんせい]{キャッシュ一貫性}[cache coherence]と呼ぶ．

\begin{figure}[htbp]
  \centering
  % Cache coherence: a write by one CPU makes the cached copies of the same
% memory location on the other CPUs stale; the hardware must update or
% invalidate them.
% Usage: % Cache coherence: a write by one CPU makes the cached copies of the same
% memory location on the other CPUs stale; the hardware must update or
% invalidate them.
% Usage: % Cache coherence: a write by one CPU makes the cached copies of the same
% memory location on the other CPUs stale; the hardware must update or
% invalidate them.
% Usage: \input{figures/l10-coherence}   (width ~116 mm)
\begin{tikzpicture}[x=1mm, y=1mm, font=\footnotesize\sffamily,
  cpu/.style={draw, fill=white, minimum width=30mm, minimum height=6mm, inner sep=1pt},
  cache/.style={draw, fill=ln-box, minimum width=30mm, minimum height=6mm,
                inner sep=1pt, font=\scriptsize\sffamily},
  lab/.style={font=\scriptsize\sffamily, align=center},
  >={Stealth[length=1.8mm]}]
  \node[cpu] (c1) at (17,14) {CPU 1: \texttt{test\_and\_set}};
  \node[cpu] (c2) at (58,14) {CPU 2};
  \node[cpu] (c3) at (99,14) {CPU 3};
  \node[cache, anchor=north, fill=white, thick] (k1) at (c1.south)
    {\texttt{lock = BUSY}};
  \node[cache, anchor=north] (k2) at (c2.south)
    {\texttt{lock = FREE}\ \iflangja{（古い）}{(stale)}};
  \node[cache, anchor=north] (k3) at (c3.south)
    {\texttt{lock = FREE}\ \iflangja{（古い）}{(stale)}};
  \draw[line width=1.6pt] (0,-6) -- (116,-6);
  \foreach \k in {1,2,3} { \draw (k\k.south) -- (k\k.south |- 0,-6); }
  \node[draw, fill=ln-box, minimum width=40mm, minimum height=6mm] (m) at (58,-15)
    {\iflangja{メモリ}{memory}: \texttt{lock = BUSY}};
  \draw (m.north) -- (58,-6);
  % write path
  \draw[->, thick, ln-accent] ([xshift=-6mm]k1.south) -- ++(0,-12) |- (m.west)
    node[lab, pos=0.25, left, text=ln-accent] {\iflangja{(1) 書き込み}{(1) write}};
  % coherence
  \draw[->, thick, ln-caution, densely dashed] (m.east) -| ([xshift=6mm]k3.south);
  \node[lab, text=ln-caution] at (89,-9.5)
    {\iflangja{(2) 更新または無効化}{(2) update or invalidate}};
  \draw[->, thick, ln-caution, densely dashed] ([xshift=8mm]m.north) -- ([xshift=8mm]k2.south);
\end{tikzpicture}
   (width ~116 mm)
\begin{tikzpicture}[x=1mm, y=1mm, font=\footnotesize\sffamily,
  cpu/.style={draw, fill=white, minimum width=30mm, minimum height=6mm, inner sep=1pt},
  cache/.style={draw, fill=ln-box, minimum width=30mm, minimum height=6mm,
                inner sep=1pt, font=\scriptsize\sffamily},
  lab/.style={font=\scriptsize\sffamily, align=center},
  >={Stealth[length=1.8mm]}]
  \node[cpu] (c1) at (17,14) {CPU 1: \texttt{test\_and\_set}};
  \node[cpu] (c2) at (58,14) {CPU 2};
  \node[cpu] (c3) at (99,14) {CPU 3};
  \node[cache, anchor=north, fill=white, thick] (k1) at (c1.south)
    {\texttt{lock = BUSY}};
  \node[cache, anchor=north] (k2) at (c2.south)
    {\texttt{lock = FREE}\ \iflangja{（古い）}{(stale)}};
  \node[cache, anchor=north] (k3) at (c3.south)
    {\texttt{lock = FREE}\ \iflangja{（古い）}{(stale)}};
  \draw[line width=1.6pt] (0,-6) -- (116,-6);
  \foreach \k in {1,2,3} { \draw (k\k.south) -- (k\k.south |- 0,-6); }
  \node[draw, fill=ln-box, minimum width=40mm, minimum height=6mm] (m) at (58,-15)
    {\iflangja{メモリ}{memory}: \texttt{lock = BUSY}};
  \draw (m.north) -- (58,-6);
  % write path
  \draw[->, thick, ln-accent] ([xshift=-6mm]k1.south) -- ++(0,-12) |- (m.west)
    node[lab, pos=0.25, left, text=ln-accent] {\iflangja{(1) 書き込み}{(1) write}};
  % coherence
  \draw[->, thick, ln-caution, densely dashed] (m.east) -| ([xshift=6mm]k3.south);
  \node[lab, text=ln-caution] at (89,-9.5)
    {\iflangja{(2) 更新または無効化}{(2) update or invalidate}};
  \draw[->, thick, ln-caution, densely dashed] ([xshift=8mm]m.north) -- ([xshift=8mm]k2.south);
\end{tikzpicture}
   (width ~116 mm)
\begin{tikzpicture}[x=1mm, y=1mm, font=\footnotesize\sffamily,
  cpu/.style={draw, fill=white, minimum width=30mm, minimum height=6mm, inner sep=1pt},
  cache/.style={draw, fill=ln-box, minimum width=30mm, minimum height=6mm,
                inner sep=1pt, font=\scriptsize\sffamily},
  lab/.style={font=\scriptsize\sffamily, align=center},
  >={Stealth[length=1.8mm]}]
  \node[cpu] (c1) at (17,14) {CPU 1: \texttt{test\_and\_set}};
  \node[cpu] (c2) at (58,14) {CPU 2};
  \node[cpu] (c3) at (99,14) {CPU 3};
  \node[cache, anchor=north, fill=white, thick] (k1) at (c1.south)
    {\texttt{lock = BUSY}};
  \node[cache, anchor=north] (k2) at (c2.south)
    {\texttt{lock = FREE}\ \iflangja{（古い）}{(stale)}};
  \node[cache, anchor=north] (k3) at (c3.south)
    {\texttt{lock = FREE}\ \iflangja{（古い）}{(stale)}};
  \draw[line width=1.6pt] (0,-6) -- (116,-6);
  \foreach \k in {1,2,3} { \draw (k\k.south) -- (k\k.south |- 0,-6); }
  \node[draw, fill=ln-box, minimum width=40mm, minimum height=6mm] (m) at (58,-15)
    {\iflangja{メモリ}{memory}: \texttt{lock = BUSY}};
  \draw (m.north) -- (58,-6);
  % write path
  \draw[->, thick, ln-accent] ([xshift=-6mm]k1.south) -- ++(0,-12) |- (m.west)
    node[lab, pos=0.25, left, text=ln-accent] {\iflangja{(1) 書き込み}{(1) write}};
  % coherence
  \draw[->, thick, ln-caution, densely dashed] (m.east) -| ([xshift=6mm]k3.south);
  \node[lab, text=ln-caution] at (89,-9.5)
    {\iflangja{(2) 更新または無効化}{(2) update or invalidate}};
  \draw[->, thick, ln-caution, densely dashed] ([xshift=8mm]m.north) -- ([xshift=8mm]k2.south);
\end{tikzpicture}

  \caption{CPU 1 が \texttt{lock} を \texttt{BUSY} に設定する．CPU 2 と
    CPU 3 のキャッシュにある複製は古くなる．キャッシュ一貫性のある
    プラットフォームでは，ハードウェアがそれらを更新または無効化する．}
  \label{fig:l10-coherence}
\end{figure}

キャッシュ一貫性を誰が担うかはプラットフォームによって異なる．キャッシュ
一貫性を持たない（NCC）プラットフォームでは，ハードウェアは書き込みを他の
キャッシュの複製へ伝播せず，古い複製への対処はソフトウェアの仕事となる．
キャッシュ一貫性のある（CC）プラットフォームでは，ハードウェアが2つの戦略
のいずれかに従ってキャッシュを一貫に保つ．
\begin{description}
  \item[書き込み無効化]
    \term[かきこみむこうか]{書き込み無効化}[write-invalidate]戦略では，
    書き込みが，他のキャッシュにあるその位置の複製を無効化する．バスや
    インターコネクトを渡るのは無効化だけなので，書き込み当たりの帯域は
    小さく，同じ CPU による以後の書き込みにはコストがかからない．ただし
    他の CPU での次の読み出しはキャッシュミスとなり，その位置を取得し直さ
    なければならない．
  \item[書き込み更新] \term[かきこみこうしん]{書き込み更新}[write-update]%
    戦略では，書き込みが新しい値を，複製を持つすべてのキャッシュへ運ぶ．
    他の CPU のスレッドは新しい値をただちに自分のキャッシュから読めるが，
    書き込み当たりのトラフィックは増える．
\end{description}
2つのうちでは書き込み無効化のほうが一般的である．いずれにせよ，一貫性機構
のコストは，共有された位置に書き込むソフトウェアが負担する．

\subsection{マルチプロセッサ上のアトミック命令}

マルチプロセッサ上では，アトミック命令はすべての CPU にわたってアトミック
でなければならない．ある CPU 上でのその実行は，他のあらゆる CPU 上での同じ
位置に対する実行を排除しなければならない．したがってここで用いるモデルでは，
アトミック命令は私的なキャッシュ複製に対しては働かない．キャッシュを迂回して
メモリ位置そのものに適用され，すべてのキャッシュを一貫に保つ．\margin{実際
のプロセッサは代わりにキャッシュラインの排他的な所有権を取得するが，コスト
は同程度である．}その結果，マルチプロセッサ上のアトミック命令は，単一 CPU
のシステム上の同じ命令よりも，また同じデータに対する通常の命令よりも高価に
なる．メモリに到達し，バスやインターコネクトを他の CPU と奪い合い，一貫性
機構を起動するからである．

スピンロックはこれをさらに悪化させる．\texttt{test\_and\_set} は，
\texttt{BUSY} を \texttt{BUSY} で置き換えるだけのときでも必ず書き込みを
行い，どの書き込みもキャッシュ一貫性機構を起動する．スピンするスレッドは
\texttt{test\_and\_set} を絶え間なく実行するので，複数のスレッドが待って
いるスピンロックは，ロック語の複製を持つすべての CPU に対して一貫性更新の
絶え間ない流れを生む．

\begin{keypoint}
  あらゆる同期機構は，テストアンドセットやコンペアアンドスワップといった
  アトミック命令の上に成り立つ．共有メモリ型マルチプロセッサでは，それらは
  すべての CPU にわたってアトミックでなければならず，ハードウェアは他の
  複製を無効化するか更新することでキャッシュを一貫に保つ．このため，
  アトミック命令は，とりわけ競合するロック語への繰り返しの書き込みは，
  高価になる．
\end{keypoint}

\section{スピンロックの性能}
\label{sec:l10-spin-perf}

マルチプロセッサ上でのアトミック命令のコストは，スピンロックをどう書くべき
かに影響する．これは Anderson \cite{anderson1990} が研究した問題である．%
\source{P3L4，スピンロックの性能指標，テストアンドセットスピンロック，
スピンロックの遅延の代替案，キューイングロック，スピンロックの性能比較；
Anderson \cite{anderson1990}．}スピンロックの実装は3つの指標で特徴づけ
られる．
\begin{description}
  \item[レイテンシ] 空いているロックをスレッドが獲得するのに要する時間．
    理想的には，スレッドはただちにアトミック命令を実行する．
  \item[遅延] ロックの解放から，スピンしているスレッドがそれを獲得するまで
    の時間．理想的には，スピンしているスレッドは解放にただちに気づく．その
    ためにはロックを絶え間なく検査する必要がある．
  \item[競合] \term[ろっくきょうごう]{ロック競合}[lock contention]，すなわち
    複数のスレッドが同じロックを同時に獲得しようとする状況が引き起こす
    オーバーヘッド．とりわけ，それらのアトミック命令が生むキャッシュ一貫性
    トラフィックである．理想的には，ロックが保持されている間，スピンして
    いるスレッドは何も書き込まない．
\end{description}
3つの指標すべてで理想的な実装は存在しない．最初の2つは待っているスレッドに
できるだけ頻繁にロックに触れることを求め，3つ目はできるだけまれにしか触れ
ないことを求めるからである．\cref{sec:l10-hw}の \texttt{test\_and\_set}
ループはレイテンシと遅延では理想的だが，反復のたびに書き込むので競合では
劣る．

\subsection{読み出しでのスピン}

\term[てすとあんどてすとあんどせっと]{テストアンドテストアンドセット}[test-and-test-and-set]%
という技法は，まず通常の命令でロックを読み，ロックが空いているように見えた
ときにだけ \texttt{test\_and\_set} を実行する．
\needspace{5\baselineskip}
\begin{lstlisting}[language=C, numbers=none]
void spinlock_lock(int *lock) {
  while (*lock == BUSY || test_and_set(lock) == BUSY)
    ;   /* スピン */
}
\end{lstlisting}
短絡評価される \texttt{||} は，読み出しが \texttt{FREE} を見つけたときに
だけ \texttt{test\_and\_set} を評価する．読み出しだけでは何も決まらない
---読み出しと \texttt{test\_and\_set} の間に別のスレッドがロックを獲得
しうる---が，その場合はアトミックな \texttt{test\_and\_set} が
\texttt{BUSY} を返してループが続くので，ロックは正しいままである．代価は，
すべての \texttt{test\_and\_set} の前に1回余分な読み出しが入ることであり，
レイテンシと遅延がわずかに増える．

この技法が競合に対して何をもたらすかは，プラットフォームによって異なる．
\begin{description}
  \item[NCC] 読み出しもアトミック命令もキャッシュからは供給できないので，
    スピンするすべてのスレッドのすべての反復がメモリに到達し，競合は単純な
    ループと変わらない．
  \item[CC，書き込み更新] ロックが保持されている間，スピンするスレッドは
    自分のキャッシュ複製を読み，トラフィックをまったく生まない．解放は
    すべてのキャッシュへ伝播されるので，待っているスレッドはすべて同じ
    瞬間に \texttt{FREE} を見て，\texttt{test\_and\_set} 命令をいっせいに
    発行する．
  \item[CC，書き込み無効化] 競合は単純なループより悪くなる．失敗するものも
    含めてどの \texttt{test\_and\_set} も他のすべてのキャッシュの複製を
    無効化するので，スピンする各スレッドは次の読み出しでミスし，ロック語を
    メモリから取得し直す．したがって1回の解放が，アトミック命令のバーストと
    それに続くキャッシュミスのバーストを引き起こす．
\end{description}
読み出しでのスピンが取り除くのは待機期間のトラフィックであって，解放の
トラフィックではない．クリティカルセクションが短く，多数のスレッドが待って
いると，解放が次々と続くのでこれらのバーストが支配的になる．書き込み無効化
のハードウェア上の激しい競合のもとで，この技法がすべての中で最も性能が悪い
のはこのためである．

\begin{example}
  読み出しでのスピンが最もよく働く軽負荷の場合を考える．4つの CPU があり，
  CPU~1 のスレッドが \SI{1}{\milli\second} ロックを保持する間，CPU 2--4 の
  スレッドがスピンする．単純な \texttt{test\_and\_set} ループで1回の反復に
  \SI{50}{\nano\second} かかるとすると，スピンする各スレッドは
  $\SI{1}{\milli\second}/\SI{50}{\nano\second} = \num{20000}$ 回反復するので，
  この保持は $3\times\num{20000} = \num{60000}$ 回のロック語への書き込みを
  要し，その1回ごとに他のキャッシュの複製を無効化または更新する．テスト
  アンドテストアンドセットでは，この保持に書き込みは1回もかからない．反復は
  読み出しであり，読み出しのほうが安価なので回数はむしろ増える．そして解放
  にかかるのは，待っているスレッド1つにつき1回，およそ3回の
  \texttt{test\_and\_set} 命令であり，そのうち1つが成功する．数百ナノ秒の
  クリティカルセクションを伴う高負荷のもとでは，そのバーストとそれに続く
  キャッシュミスが解放のたびに繰り返されることになる．
\end{example}

\subsection{遅延を用いたスピンロック}

バーストが生じるのは，待っているスレッドがすべて同じ瞬間に解放へ反応する
からである．ロックに再び触れる前に少し待たせれば，反応がばらける．遅延は，
ロックが空いたのを見つけた後に置くことができる．
\needspace{9\baselineskip}
\begin{lstlisting}[language=C, numbers=none]
void spinlock_lock(int *lock) {
  while (*lock == BUSY || test_and_set(lock) == BUSY) {
    while (*lock == BUSY)
      ;        /* 保持中は読み出しでスピン */
    delay();   /* 空いた: 少し待つ */
  }
}
\end{lstlisting}
ロックの獲得に失敗したスレッドは，これまでどおり読み出しでスピンし，
ロックが空いたらアトミック命令を試みる前に待つ．そのため待っているスレッド
は互いに異なる時刻に \texttt{test\_and\_set} を試み，1回の解放に反応する
スレッドの数が減る．空いているロックに出会ったスレッドはループの本体に入ら
ないので，レイテンシは読み出しでのスピンと同じである．競合は著しく改善する
が，遅延は悪化する．ロックがすでに空いているのに，スレッドが自分の遅延の
中に留まっていることがあるからである．

遅延は，反復のたびに置くこともできる．
\begin{lstlisting}[language=C, numbers=none]
void spinlock_lock(int *lock) {
  while (*lock == BUSY || test_and_set(lock) == BUSY)
    delay();
}
\end{lstlisting}
待っているスレッドは，絶え間なくではなく遅延の期間ごとに1回ロックを参照
するようになる．これは一貫性トラフィックも減らし，読み出しでのスピンと
違って，すべての参照がメモリに到達する NCC プラットフォームでも効果がある．
代価はやはり遅延である．解放が，最長で1遅延期間のあいだ気づかれないままに
なりうる．

どちらの変種も遅延の値を必要とする．\emph{静的な}遅延は，各スレッドを区別
する何か，例えば CPU 識別子から導いた固定値を各スレッドに与え，試行が
仕組みとしてばらけるようにする．単純だが，他に待っているスレッドがない
ときでもスレッドを遅らせる．\emph{動的な}遅延は，スレッドが感じ取る競合
---例えば失敗した \texttt{test\_and\_set} の試行回数から推定する---と
ともに広がる範囲から無作為に選ばれる．\caution{試行の失敗は，待っている
スレッドが多いことではなく，クリティカルセクションが長いことを意味して
いるかもしれない．}

\subsection{キューイングロック}

遅延は，解放への反応を時間的にばらけさせる．キューイングロックは，ちょうど
1つのスレッドだけを反応させることで，バーストそのものを避ける．

\begin{definition}[キューイングロック]
  \term[きゅーいんぐろっく]{キューイングロック}[queueing lock]は，ロックを
  要求したすべてのスレッドに，キュー内の場所と専用のフラグを与える．待って
  いるスレッドは自分のフラグ上でスピンし，ロックを解放するスレッドはキュー
  内の後継のフラグを設定する．これによって後継が所有者となる．
\end{definition}

キューはスレッドごとに1つのフラグからなる配列であり，各フラグは
\texttt{HAS\_LOCK} か \texttt{MUST\_WAIT} のいずれかを取る．これに次の空き
場所の添字が伴い，到着したスレッドは \texttt{read\_and\_increment} でそれを
得る．
\begin{lstlisting}[language=C, numbers=none]
int flags[N];   /* flags[0] = HAS_LOCK，残りは待機 */
int queuelast = 0;

int queue_lock(void) {
  int me = read_and_increment(&queuelast) % N;
  while (flags[me] == MUST_WAIT)
    ;                   /* 専用のフラグ上でスピン */
  flags[me] = MUST_WAIT;  /* 次の周回に備えて戻す */
  return me;
}

void queue_unlock(int me) {
  flags[(me + 1) % N] = HAS_LOCK;
}
\end{lstlisting}
解放は1つのフラグだけを書き込むので，起こされる待機スレッドはちょうど1つ
であり，無効化または更新されるキャッシュ複製も1つだけである．3つの指標に
照らすと，キューイングロックは次のようになる．
\begin{itemize}
  \item \emph{レイテンシ}は悪い．\texttt{read\_and\_increment} は
    \texttt{test\_and\_set} より複雑なアトミック命令であり，獲得のたびに
    それを実行してから配列を添字で引くからである．
  \item \emph{遅延}は良い．所有者がロックを後継へ直接手渡すので，後継が
    自分で解放を検出する必要がない．
  \item \emph{競合}ははるかに良い．異なるスレッドのフラグが異なるキャッシュ
    ラインに置かれている限り，1回の解放が生むトラフィックは1つのフラグ分
    だけである．
\end{itemize}
その代わりにロックは，1語ではなく $N$ 項目の配列を占める．ここで $N$ は，
そのロックを要求しうるスレッドの数である．また \texttt{read\_and\_increment}
はすべてのプラットフォームで利用できるわけではない．

\subsection{比較}

\cref{tab:l10-spin-compare}は評価をまとめたものである．Anderson は，
書き込み無効化キャッシュを持つ 20 CPU の Sequent Symmetry 上でこれらの実装
を測定した \cite{anderson1990}．高負荷---多数の CPU が同じロックを奪い合う
---のもとではキューイングロックが最も良く，テストアンドテストアンドセットが
最も悪く，遅延を用いたロックはその中間であった．軽負荷では順序が逆転した．
テストアンドテストアンドセットはレイテンシと遅延が最も小さいので最も良く，
キューイングロックは獲得のたびに帳簿付けの代価を払うので最も悪い．すべての
条件で最良の実装は存在しない．選択は，予想される競合と，プラットフォームが
提供する命令に依存する．

\begin{table}[htbp]
  \centering
  \caption{3つの指標に照らしたスピンロックの実装．各項目は互いに相対的な
    もので，測定された順位は負荷に依存する．}
  \label{tab:l10-spin-compare}
  \small
  \begin{tabular}{@{}llll@{}}
    \toprule
    実装 & レイテンシ & 遅延 & 競合 \\
    \midrule
    \texttt{test\_and\_set}        & 良 & 良 & 悪 \\
    テストアンドテストアンドセット & 中 & 中 & 書き込み無効化では悪 \\
    解放後に遅延                   & 中 & 悪 & 改善 \\
    参照のたびに遅延               & 中 & 悪 & 改善，NCC でも有効 \\
    キューイングロック             & 悪 & 良 & 最良 \\
    \bottomrule
  \end{tabular}
\end{table}

\begin{keypoint}[章のまとめ]
  ミューテックスと条件変数は誤りやすく，表現力に限界があり，ハードウェアの
  支援に依存する．スピンロックはブロックする代わりにビジーウェイトし，
  複数 CPU 上の短いクリティカルセクションでのみ引き合う．セマフォは与えられ
  た数までのスレッドを受け入れ，バイナリセマフォはミューテックスとして働く．
  リーダ・ライタロックは共有モードと排他モードを提供し，モニタはエントリ
  手続きの中でロックを暗黙的に獲得・解放する．シリアライザ，パス式，バリア，
  ランデブーポイント，および RCU のようなウェイトフリーの方式がこれを補う．
  そのすべてが，テストアンドセットやコンペアアンドスワップといったアトミック
  命令に依存しており，共有メモリ型マルチプロセッサではさらに CPU のキャッシュ
  を一貫に保たなければならない．テストアンドセットの前にロックを検査する
  ことは軽負荷では有効だが，書き込み無効化のプラットフォーム上の高負荷では
  すべての中で最も性能が悪い．待っているスレッドを遅延させること，とりわけ
  専用のフラグ上でキューに並ばせることが，競合するロックの一貫性トラフィック
  を抑える．
\end{keypoint}

\end{document}
